\section{Privacy–Security Frontier：用工程扩大可选空间}

Palantir 从反恐场景起步，天然面对隐私和公民自由争议。Sankar 把政治与工程的分工表述为：政治决定社会愿意位于 privacy/security trade-off 的哪个点，工程则应把 efficient frontier 向外推，使同等隐私保护下获得更多安全，或同等安全能力下减少数据侵入。这个框架有用，但“向外推”必须由可测量机制支持，而不能成为扩大监控的修辞。

\subsection{Efficient frontier 的含义}

\term{Efficient frontier} 是无法在不牺牲另一目标的情况下继续改善某一目标的边界。设隐私保护为 $P$，安全能力为 $S$，工程改进希望让可行集合从 $\mathcal{F}_0$ 扩展到 $\mathcal{F}_1$，使某些点满足相同 $P$ 下更高 $S$。这里每个符号都需要实际指标，例如数据最小化、误报率、案件处理时间和未经授权访问，而不是抽象“更多安全”。

\begin{figure}[H]
\centering
\includegraphics[width=0.94\textwidth]{images/12-privacy-security-frontier.png}
\caption{工程改进可以推动 privacy–security capability frontier，而社会仍需选择可接受点。}
\figsource{依据字幕 39:30--42:00 重绘，`images/12-privacy-security-frontier.png`。}
\end{figure}

\begin{knowledgebox}{读图：哪些机制可能真正推动 frontier}
Fine-grained access control 减少无关人员看到数据；purpose limitation 限制用途；immutable audit 让滥用可追责；privacy-preserving linkage 在不复制全部原始数据时完成匹配；retention policy 自动删除过期数据。每种机制都应接受 red-team、误用测试和独立监督。
\end{knowledgebox}

\begin{warningbox}{Capability 不能替代正当程序}
即使系统技术上能更精确地筛选对象，也仍需回答谁授权、谁可申诉、错误怎样纠正、数据保留多久、监督者能看到什么。Frontier 改善扩大了制度选择空间，不会自动选出合法或正义的使用方式。
\end{warningbox}

\subsection{本章小结}

Privacy 与 security 不是只能靠口号争夺的单一滑块。Access control、数据最小化、审计和隐私保护计算可以降低部分冲突，但最终目标和授权仍由政治与法律决定。技术团队的责任是让 trade-off 可测量、可限制和可复查。

